\documentclass[10pt]{article}
\usepackage[margin=0.85in]{geometry}
\usepackage{amsmath,amssymb}
\usepackage[hidelinks]{hyperref}

\begin{document}
\begin{center}
  {\Large\bfseries The fundamental group of a \(2\)-knot complement\par}
  \smallskip
  Disproof of Conjecture 00000000656
\end{center}

Let \(K\) be an embedded \(2\)-sphere in the interior of \(B^4\),
as for a closed \(2\)-knot or a spun knot placed in a ball. All
homology and cohomology groups below have integer coefficients.
We show that the fundamental group of \(B^4\setminus K\) surjects
onto \(\mathbb Z\).

\section*{The capped complement}

Glue a second \(4\)-ball to \(B^4\) along their boundary to view
\(K\) as a subset of \(S^4\). Set \(X=S^4\setminus K\).
Alexander duality applies because \(K\) is compact, nonempty,
locally contractible, and proper in \(S^4\)
\cite[Corollary~3.45]{hatcher}. It gives
\[
  \widetilde H_1(X)\cong \widetilde H^2(K)\cong\mathbb Z,
  \qquad
  \widetilde H_0(X)\cong \widetilde H^3(K)=0 .
\]
Here \(K\cong S^2\) has a CW structure with one \(0\)-cell and
one \(2\)-cell. In degree two its cellular cochain group is
\(\mathbb Z\), with neither incoming coboundaries nor outgoing
differential. Its degree-three cochain group is zero.

\section*{Removing the cap}

Choose an open collar \(A\) of the gluing \(S^3\) in \(S^4\)
disjoint from \(K\). Cover \(X\) by the open sets \(U\) and \(V\)
obtained by adjoining \(A\) to the respective ball interiors, with
\(K\) removed from the first. Their intersection is \(A\), which
retracts onto \(S^3\). Collapsing the extra half-collars gives
deformation retractions of \(U\) to \(B^4\setminus K\) and \(V\)
to the cap ball. In the reduced Mayer--Vietoris sequence the segment
\[
  0=\widetilde H_1(A)\ \longrightarrow\
  \widetilde H_1(U)\oplus\widetilde H_1(V)\
  \longrightarrow\ \widetilde H_1(X)\
  \longrightarrow\ \widetilde H_0(A)=0
\]
shows that \(H_1(B^4\setminus K)\cong H_1(X)\cong\mathbb Z\),
since \(V\) is contractible. The segment from
\(\widetilde H_0(A)=0\) to \(\widetilde H_0(X)=0\)
also gives \(\widetilde H_0(U)=0\), so \(B^4\setminus K\)
is path connected.

Degree-one Hurewicz identifies \(H_1\) of a path-connected space
with the abelianization of its fundamental group
\cite[Theorem~2A.1]{hatcher}. Therefore
\[
   \pi_1(B^4\setminus K)\twoheadrightarrow \mathbb Z .
\]
The group is nontrivial for every such \(2\)-knot, including every
spun knot. This rules out the proposed family.

\section*{Lean verification}

The Lean~4.19.0 project computes the degree-two cellular cohomology
of the two-cell sphere as cocycles modulo coboundaries. It proves
evaluation on the \(2\)-cell has a two-sided inverse and preserves
addition.
Given the maps and exactness from the topological results above,
it derives surjectivity through the Mayer--Vietoris cap step and
constructs distinct fundamental-group elements over the classes
of \(0\) and \(1\). Run \verb|lake build| in \verb|lean/|
to verify the formal argument.

\begin{thebibliography}{1}
\bibitem{hatcher}
A.~Hatcher, \emph{Algebraic Topology}, Cambridge University Press, 2002,
\href{https://pi.math.cornell.edu/~hatcher/AT/AT.pdf}{author's online edition}.
\end{thebibliography}

\end{document}
